\documentclass[12pt,a4paper]{article}
\usepackage[margin=24mm,top=20mm,bottom=18mm]{geometry}
\usepackage{fontspec}
\setmainfont{Cambria}\setmonofont{Consolas}
\usepackage{polyglossia}
\setdefaultlanguage{russian}
\usepackage[colorlinks=true,urlcolor=blue]{hyperref}
\hypersetup{pdftitle={offix — описание сервиса},pdfauthor={Иван Бычков и Роман Емельянов}}
\pagestyle{empty}
\setlength{\parindent}{0pt}
\setlength{\parskip}{7pt}
\linespread{1.10}
\newcommand{\topic}[1]{\par\vspace{4pt}{\large\bfseries #1}\par\nopagebreak}
\begin{document}
{\LARGE\bfseries offix}\par
{\large Описание сервиса для сотрудников университета}\par
\vspace{2pt}
Иван Бычков и Роман Емельянов разрабатывают offix для конкурса «Академия Финтеха» Финансового университета и ВТБ в рамках трека «Эффективность~360». Сервис предназначен для совместной работы подразделений. В нём можно поставить общую цель, распределить задачи между сотрудниками и отслеживать результат по выбранному показателю.

\topic{Цели и задачи}
В карточке цели руководитель указывает, что должно измениться, кто отвечает за результат и к какому сроку его нужно получить. Для оценки результата выбирают показатель, фиксируют его исходное значение и задают целевое. В университете таким показателем может быть среднее время обработки обращений или срок согласования документов.

Работу над целью делят на задачи. У каждой есть исполнитель и срок; при необходимости назначают согласующего. Задачи могут идти параллельно или зависеть друг от друга. На карте цели видно, какой отдел ждёт результата от коллег и задержка каких шагов влияет на общий срок. Если сотрудник не может продолжить работу, он указывает причину. Руководитель видит, какие задачи это затрагивает, и принимает решение о переносе сроков.

Ход работы и результат учитываются отдельно. Статусы задач показывают, что уже сделано. Замеры показателя позволяют увидеть, насколько организация приблизилась к цели. Их можно вносить вручную или загружать из таблицы.

\topic{Работа сотрудников и руководителей}
\textbf{Текущие задачи.} Сотрудник видит свои поручения и сроки, отмечает выполненные шаги и передаёт результат на согласование. Вопрос можно обсудить прямо в задаче или в чате с коллегами.

\textbf{Обзор и отчёты.} На экране «Пульс» руководитель видит цели подразделений и причины риска по каждой из них. Для совещания можно подготовить сводку по показателям и срокам или настроить получение отчёта по электронной почте.

\textbf{База знаний.} Сотрудники могут хранить регламенты, инструкции и образцы документов. Материалы, относящиеся к определённой цели, доступны из её карточки.

\vspace{7pt}\hrule\vspace{6pt}
\textbf{Демонстрационная версия}\quad\url{http://next.147-45-43-34.sslip.io/login}\par
На странице входа выберите сотрудника демонстрационного университета. Под учётной записью ректора доступны общий обзор, цели и карты задач.
\end{document}


